
\documentclass[12pt]{article}

% Chinese support
\usepackage{xeCJK}
\setCJKmainfont{Microsoft YaHei}
\setCJKsansfont{Microsoft YaHei}
\setCJKmonofont{FangSong}

% Sans-serif font (fontspec for XeLaTeX)
\usepackage{fontspec}
\setmainfont{Arial}
\setsansfont{Arial}
\setmonofont{Courier New}

% Sans-serif math font
\usepackage[T1]{fontenc}
\usepackage{sfmath}
\usepackage{amssymb}

\newcommand{\AF}{\mathrm{AF}}

% Page margins
\usepackage[margin=0.1cm]{geometry}

% Section title sizes
\usepackage{titlesec}
\setcounter{secnumdepth}{4}
\titleformat{\section}{\fontsize{36}{43}\bfseries}{\thesection}{1em}{}
\titleformat{\subsection}{\fontsize{28}{34}\bfseries}{\thesubsection}{1em}{}
\titleformat{\subsubsection}{\fontsize{22}{27}\bfseries}{\thesubsubsection}{1em}{}
\titleformat{\paragraph}[runin]{\sffamily\fontsize{22}{27}\bfseries}{\theparagraph}{1em}{}

\begin{document}
\sffamily
\huge

\section{crossing}
\paragraph{}给定一个谱序列E，收敛到的A
\paragraph{}假设谱序列都从至少第0页或正的某页开始
\paragraph{}假设E的微分满足：$d_n$将filtration增加n

\paragraph{}essential 微分指作为$E_n$页上的映射，打过去非零

\subsection{定义2.7（来自2412.10879）}

\begin{enumerate}
    \item For $x\in V$，假设x的filtration为s, we say that the an 微分 $d_n(x)=y$ has a crossing that hits filtration $p$ for some $p \leq s+n$, if there exists an element $x^\prime\in V$，在filtraton a+s， with $a>0$ and an essential 微分 $d_m(x^\prime)=y^\prime$ for $0\neq y^\prime$，在filtration s+a+m, such that
    $$p = \textup{Fil}(y^\prime) = s+a+m\le \textup{Fil}(y) = s+n.$$
    \item We say that an 微分 $d_n(x)=y$ has no crossing that hits the range Fil $\ge p$ if there does not exist an element $x^\prime\in V$，filtration $s+a$ with $a>0$ and an essential 微分 $d_m(x^\prime)=y^\prime$ for $0\neq y^\prime\in V$，filtration s+a+m， such that
    $$p\le \textup{Fil}(y^\prime) = s+a+m\le \textup{Fil}(y) = s+n.$$
    \item We say that an 微分 $d_n(x)=y$ has no crossing if it has no crossing that hits the range Fil $\ge Fil(x) + 1$.
\end{enumerate}



\end{document}
